\documentclass[10pt,conference,twocolumn,a4paper]{article}

\usepackage[margin=1.8cm]{geometry}
\usepackage[utf8]{inputenc}
\usepackage[T1]{fontenc}
\usepackage{microtype}
\usepackage{amsmath,amssymb}
\usepackage{booktabs}
\usepackage{graphicx}
\usepackage{xcolor}
\usepackage{hyperref}
\usepackage{enumitem}
\usepackage{titlesec}

\hypersetup{colorlinks=true, linkcolor=blue!50!black, citecolor=blue!50!black,
            urlcolor=blue!50!black}

\titlespacing*{\section}{0pt}{0.7\baselineskip}{0.25\baselineskip}
\titlespacing*{\subsection}{0pt}{0.5\baselineskip}{0.15\baselineskip}
\setlist{nosep,leftmargin=1.1em}
\setlength{\columnsep}{0.7cm}
\setlength{\parskip}{2pt}
\setlength{\abovecaptionskip}{2pt}
\setlength{\belowcaptionskip}{0pt}

\title{\vspace{-2em}\textbf{P1 -- Hierarchical Concept-Based Explainable-by-Design Models:\\
A Two-Backbone Study on CUB-200-2011}}
\author{Sina Aghakhani \\
\small Explainable and Trustworthy AI -- Politecnico di Torino, A.Y.\ 2025/2026 \\
\small Reference teachers: Eliana Pastor, Eleonora Poeta}
\date{\vspace{-1.5em}}

\begin{document}
\maketitle
\thispagestyle{empty}

\begin{abstract}\small
We address project P1 by designing a Hierarchical Concept Bottleneck Model
(H-CBM) in which the 312 fine attributes of CUB-200-2011 are not a flat set
but are explicitly organised under 13 coarse body parts. Hierarchy is enforced
\emph{architecturally} through a multiplicative parent-conditioned gate that
is identical at training and inference, removing the need for soft consistency
penalties. We instantiate the same H-CBM on two backbones (ResNet-50 and
ViT-B/16) and run an identical interpretability suite on both: per-attribute
AUC, class-prototype oracle, test-time intervention (TTI), per-part causal
necessity, calibration, classifier-weight structure and pairwise agreement.
ViT reaches 86.56\% top-1 vs.\ 61.94\% for ResNet, with smaller
class-prototype oracle gap (13.4 vs.\ 35.0~pp), better calibration
(ECE 0.126 vs.\ 0.180) and a strictly monotone TTI curve
($61.9{\to}96.9\%$ ResNet, $86.6{\to}100.0\%$ ViT). Per-attribute analysis
shows ViT systematically wins on shape/size and rare colours
(\emph{rufous, purple}), ResNet wins on a few common colours
(\emph{pink, green eye/throat}). Cohen's $\kappa{=}0.632$ and an
oracle-selection ensemble ceiling of 89.21\% indicate genuine
complementarity rather than redundancy.
\end{abstract}

\section{Project goal and deliverables}
The project brief~\cite{p1brief} requires:\
(D1) a literature review on concept-bottleneck and hierarchical
interpretable models;\
(D2) the identification of research gaps in current concept-based XAI;\
(D3) an implementation that introduces hierarchy into the prediction
pathway, with coarse concepts \emph{supporting} fine ones; and\
(D4) an evaluation that combines predictive accuracy, concept quality and
hierarchy-aware analyses. We satisfy each deliverable as follows.

\subsection{D1 -- Literature review}
We surveyed CBMs~\cite{koh2020}, hierarchical CBMs~\cite{pittino2023,sun2024},
hierarchical prototypes~\cite{hase2019}, mistake-severity
metrics~\cite{bertinetto2020} and concept-based XAI taxonomies~\cite{poeta2025},
plus the supporting tools we used: Focal Loss~\cite{lin2017} and
multi-task uncertainty weighting~\cite{kendall2018}.

\subsection{D2 -- Research gaps}
Four gaps are addressed by the design (Tab.~\ref{tab:gaps}): flat concept
sets, soft-only coarse-to-fine constraints, lack of hierarchy-aware
evaluation, and classifier leakage that bypasses the bottleneck.
\begin{table}[h]
\centering\footnotesize
\begin{tabular}{@{}p{0.07\linewidth}p{0.40\linewidth}p{0.40\linewidth}@{}}
\toprule
\textbf{Gap} & \textbf{Prior best} & \textbf{Our solution}\\
\midrule
G1 & flat concept list~\cite{koh2020} & 13/312 hierarchy parsed from attribute names; no extra labels\\
G2 & soft consistency loss~\cite{pittino2023} & multiplicative gate inside forward pass\\
G3 & top-1 only & paired suite: AUC, oracle, TTI, part-necessity, ECE, $W$-structure\\
G4 & joint training, side-channel & strict $\mathrm{Linear}(312{\to}200)$, sequential 3-phase training\\
\bottomrule
\end{tabular}
\caption{Mapping from research gaps to architectural / evaluation choices.}
\label{tab:gaps}
\end{table}

\subsection{D3 -- Implementation}
\textbf{Dataset (notebook~01).} CUB-200-2011, 11{,}788 images, 200 species,
stratified into train / val / test splits of 5{,}394 / 600 / 5{,}794. The
312 binary attributes are kept only at certainty $\geq 3$
(${\sim}259$ trusted attributes per image on average).

\noindent\textbf{Hierarchy.} The 312 attributes are deterministically parsed
by prefix into 13 body parts (\texttt{back, belly, bill, breast, crown, eye,
forehead, head, leg, nape, tail, throat, wing}) with $\{34, 34, 27, 19, 15,
14, 15, 11, 15, 15, 40, 15, 24\}$ children respectively; 34 global
attributes (\texttt{shape}, \texttt{size}, \dots) have no parent and are
masked with $1$. Part labels are the OR-aggregation of their children, so
\emph{no extra annotation is required}.

\noindent\textbf{Architecture.} Given backbone features $F$, a coarse head
produces $p_c\in[0,1]^{13}$ and a fine head $p_{f,\text{raw}}\in[0,1]^{312}$.
The hierarchy is enforced by the \emph{same} gate at train and test:
\begin{equation}
p_f[i] \;=\; \sigma(z_f[i]) \cdot \bigl(0.5 + 0.5\,p_c[\mathrm{parent}(i)]\bigr).
\label{eq:mask}
\end{equation}
The classifier is a single linear layer
$\hat{y}=\mathrm{softmax}(W p_f + b)$, $W\in\mathbb{R}^{200\times 312}$,
with no residual side-channel from $F$.

The gate in (\ref{eq:mask}) is in $[0.5,1]$ rather than $[0,1]$ on purpose:
a hard $\times p_c[\mathrm{parent}]$ collapses children to zero whenever
$p_c\!\approx\!0.5$ early in training and the fine head never recovers;
the soft form keeps gradients alive while still imposing a parent-conditioned
ceiling. Crucially, training and inference use the \emph{same} formula, so
the classifier never sees out-of-distribution $p_f$ values at test time.

\noindent\textbf{Training (3 phases).} Phase~1 trains the heads on a frozen
backbone with $\lambda_c L_{\text{coarse}} + \lambda_f L_{\text{fine}} +
0.1\,L_{\text{task}}$ (weighted BCE / Focal Loss / CE-LS). Phase~2 calibrates
\emph{only} the classifier on the model's own \emph{predicted}~$p_f$ to
match the test-time distribution. Phase~3 jointly fine-tunes everything
with concept-dominant weights ($2/2/1$) and MixUp~$\alpha{=}0.2$. The
backbone uses an order-of-magnitude lower LR than the heads; for ViT-SWAG
the backbone LR is reduced further (1e-6) to preserve attention features.

\noindent\textbf{Two backbones, identical heads.} We instantiate the same
H-CBM on ResNet-50 (336\,px input, ImageNet-V2; best at epoch 127, val
64.17\%) and ViT-B/16 (384\,px, SWAG-E2E; 86.1\,M backbone + 0.82\,M head
parameters, best at epoch 66, val 86.50\%). Everything above the backbone
-- the 13/312 hierarchy, both heads, the gate in (\ref{eq:mask}), the
linear classifier -- is byte-identical, so any observed difference is
attributable solely to the backbone choice.

\subsection{D4 -- Evaluation}
The two trained models are loaded into a shared evaluation harness
(notebook~04) which builds one paired cache and runs nine identical analyses.
Two methodological points required care.

\noindent\textbf{(i) Class-prototype oracle, not raw GT.} The classifier is
trained on continuous \emph{predicted} $p_f$. Empirically the gap between
attribute-present and attribute-absent activations is small: on ResNet,
$\mathbb{E}[p_f\mid y{=}1]{=}0.395$ vs.\ $\mathbb{E}[p_f\mid y{=}0]{=}0.287$
(gap $+0.107$ only), so most class-discriminative signal lives in the
continuous noise pattern, not in binary attribute identity. Feeding raw
$\{0,1\}$ ground truth as a fake oracle is therefore out-of-distribution and
collapses to $\approx\!11\%$ on both backbones (10.51\% R, 11.89\% V). The
honest ceiling is the class-prototype oracle
$\mathrm{proto}[k]=\mathbb{E}_{i:y_i=k}[p_f(x_i)]$, which lives exactly on
the training distribution.

\noindent\textbf{(ii) TTI by mixing toward the prototype.} Section~5 mixes
$p_f$ with $\mathrm{proto}[y]$ over fractions
$r\in\{0,0.1,0.25,0.5,0.75,1\}$ (3 random seeds per fraction): $r{=}0$
recovers actual accuracy, $r{=}1$ recovers the oracle. A monotonically
rising curve is the genuine evidence that the head listens to the concept
layer.

\noindent\textbf{(iii) Per-part necessity by zeroing child columns.} Because
the gate (\ref{eq:mask}) is in $[0.5,1]$, setting $p_c[p]{=}0$ only halves
the part's children. The honest ablation is to zero the corresponding
columns of $p_f$, which removes the branch from $W p_f$ entirely.

\section{Quantitative results}
Table~\ref{tab:headline} reports the verdict produced by notebook~04 on the
CUB test split ($N{=}5794$).

\noindent\textbf{Concept fidelity (\S3).} The 312 per-attribute AUC
distributions overlap heavily: the means are 0.636 vs.\ 0.656. Both models
share the same easy/hard pattern -- yellow body colours top the AUC ranking
on both backbones (0.87--0.91), while pink/purple colours and rare shapes
sit at the bottom (0.36--0.46) -- a direct consequence of CUB's class
imbalance on these attributes. ViT systematically wins on shape/size
(\texttt{shape:owl-like} +0.27, \texttt{size:very\_large} +0.20) and rare
colours (\texttt{bill\_color:rufous} +0.32, \texttt{under\_tail:purple}
+0.28); ResNet wins on a small set of common colours
(\texttt{eye\_color:green} $-$0.26, \texttt{throat:pink} $-$0.23). The
asymmetry is concept-localised, not uniform.

\noindent\textbf{Causal control (\S5).} TTI is strictly monotone for both
backbones (Tab.~\ref{tab:tti}), and each curve reaches its own
class-prototype oracle at $r{=}1$ by construction, confirming that the
entire decision is a faithful function of the 312-dim bottleneck.

\begin{table}[h]
\centering\footnotesize
\setlength{\tabcolsep}{6pt}
\begin{tabular}{@{}rrr@{}}
\toprule
$r$ (\% replaced) & ResNet-50 & ViT-B/16 \\
\midrule
0\,\%   & 61.94 & 86.56 \\
10\,\%  & 66.18 & 88.54 \\
25\,\%  & 72.94 & 91.27 \\
50\,\%  & 85.07 & 95.79 \\
75\,\%  & 93.48 & 98.88 \\
100\,\% & 96.91 & 100.00 \\
\bottomrule
\end{tabular}
\caption{TTI accuracy (\%) when a fraction $r$ of $p_f$ columns is replaced
by the in-class prototype, mean over 3 seeds.}
\label{tab:tti}
\end{table}

\noindent\textbf{Complementarity (\S9).} ViT is uniquely correct on $1{,}580$
images, ResNet on $154$, both wrong on $625$ -- so each model still has
recoverable errors. Cohen's $\kappa{=}0.632$ confirms the disagreements are
not trivial. ResNet's unique strengths cluster on large seabirds
(albatrosses, gulls); ViT's on small passerines (warblers, vireos,
flycatchers, sparrows), a backbone-specific signal that survives the
shared concept layer.

\begin{table}[h]
\centering\footnotesize
\setlength{\tabcolsep}{4pt}
\begin{tabular}{@{}lrrl@{}}
\toprule
\textbf{axis} & \textbf{ResNet-50} & \textbf{ViT-B/16} & \textbf{winner}\\
\midrule
top-1 acc (\%)            & 61.94  & \textbf{86.56}  & ViT \\
mean per-class acc (\%)   & 62.13  & \textbf{86.64}  & ViT \\
NLL                       & 1.459  & \textbf{0.697}  & ViT \\
mean attr-AUC             & 0.636  & \textbf{0.656}  & ViT \\
attr ECE                  & 0.180  & \textbf{0.126}  & ViT \\
attr Brier                & 0.154  & \textbf{0.112}  & ViT \\
oracle acc (\%)           & 96.91  & \textbf{100.00} & ViT \\
gap oracle$-$actual (pp)  & $+34.97$ & \textbf{$+13.44$} & ViT \\
$W$ sparsity (L1/L2)      & \textbf{14.50} & 14.72  & ResNet \\
$W$ part Herfindahl       & \textbf{0.093} & 0.092  & ResNet \\
$\sum$ part-drops (pp)    & \textbf{6.52}  & 6.45   & ResNet \\
\midrule
\multicolumn{4}{c}{axes-won score: \textbf{ResNet=4, ViT=7}}\\
\bottomrule
\end{tabular}
\caption{Final verdict, paired evaluation on the CUB test split.}
\label{tab:headline}
\end{table}

\section{Discussion and conclusion}\label{sec:disc}
\textbf{Interpretability of the bottleneck holds.} For \emph{both} models the
TTI curve is strictly monotone and reaches the prototype-oracle ceiling at
$r{=}1$, which means the entire decision is a faithful function of the 312
concepts -- there is no shortcut around the bottleneck.

\noindent\textbf{ViT's gain is in how it \emph{uses} concepts.} The mean
attribute-AUC differs only by 2~points (0.656 vs.\ 0.636), but the oracle
gap differs by 21.5~pp. Concepts are not dramatically more accurate on ViT;
ViT's classifier is much closer to its own ceiling, suggesting the
backbone shapes the $p_f$ distribution into a form on which a linear head
can separate fine-grained species. The $100.00\%$ ViT oracle is partly an
artefact of evaluating an already-strong classifier on its own in-class
average $p_f$ -- with $\sim\!30$ test images per class the prototype is a
clean class signature -- which is why we report the gap rather than the
absolute value as the headline number.

\noindent\textbf{ResNet leaves headroom on the table.} A 35-pp oracle gap
combined with monotone TTI tells us that the dominant remaining error is
in the concept extractor itself; concept-improvement (better backbones,
attention-pooled features) would translate almost one-for-one into accuracy.

\noindent\textbf{Per-part reliance is qualitatively similar.} Both models
list \texttt{bill}, \texttt{tail} and \texttt{belly} in their top-3
most-necessary parts (max drop 1.40~pp ResNet, 1.09~pp ViT;
$\sum$~drops $\approx\!6.5$~pp for both). The $W$ matrices show comparable
sparsity ($\sim\!14.5$ L1/L2) and almost identical part-Herfindahl ($\approx 0.09$),
so the species \emph{definitions} learned on top of the two backbones are
structurally close, and the accuracy gap is not explained by one model
``cheating'' through a more concentrated $W$. As an additional
sanity check, the closest pairs in concept space recovered
\emph{independently} on both backbones are taxonomically meaningful
(\textit{Common Tern}\,$\leftrightarrow$\,\textit{Forster's Tern},
\textit{California Gull}\,$\leftrightarrow$\,\textit{Herring Gull},
\textit{Loggerhead}\,$\leftrightarrow$\,\textit{Great Grey Shrike}),
showing that the bottleneck has internalised genuine ornithological
structure.

\noindent\textbf{Complementarity.} The oracle ensemble gains $+2.65$~pp
over the better single model, and the per-class only-correct lists are
species-specific (warblers, vireos, sparrows favour ViT; large gulls and
albatrosses favour ResNet), suggesting the two backbones extract
complementary inductive biases that the concept layer surfaces in a
human-readable way.

\noindent\textbf{Conclusion.} An H-CBM with an architecturally enforced
parent-conditioned gate, a strict linear bottleneck and a sequential
training schedule satisfies all four deliverables of P1: it adds explicit
hierarchy without extra labels, evaluates concept fidelity and causal
control with hierarchy-aware metrics, and remains directly comparable
across backbones. ViT-B/16 is the better backbone for this design on CUB,
but ResNet-50 produces equivalently \emph{interpretable} structure --
the bottleneck pays off for both.

\bibliographystyle{IEEEtran}
\begin{thebibliography}{9}\footnotesize
\bibitem{koh2020} P.W.\ Koh \emph{et al.}, ``Concept Bottleneck Models,'' \emph{ICML}, 2020.
\bibitem{pittino2023} F.\ Pittino \emph{et al.}, ``Hierarchical CBM for Vision,'' \emph{EAAI}, 2023.
\bibitem{sun2024} X.\ Sun \emph{et al.}, ``Supervised Hierarchical CBM,'' under review, 2024.
\bibitem{poeta2025} E.\ Poeta \emph{et al.}, ``Concept-based XAI: A Survey,'' \emph{ACM CSUR}, 2025.
\bibitem{hase2019} P.\ Hase \emph{et al.}, ``Interpretable Image Recognition with Hierarchical Prototypes,'' \emph{AAAI HCOMP}, 2019.
\bibitem{bertinetto2020} L.\ Bertinetto \emph{et al.}, ``Making Better Mistakes,'' \emph{CVPR}, 2020.
\bibitem{lin2017} T.-Y.\ Lin \emph{et al.}, ``Focal Loss for Dense Object Detection,'' \emph{ICCV}, 2017.
\bibitem{kendall2018} A.\ Kendall \emph{et al.}, ``Multi-Task Learning Using Uncertainty,'' \emph{CVPR}, 2018.
\bibitem{p1brief} ``P1 -- Hierarchical Concept-Based Explainable-by-Design Models,'' XAI Course brief, PoliTo, 2025/2026.
\end{thebibliography}

\end{document}
